\section{Evaluation Design}\label{sec:evaluation}
Since inferred MRs are not ground-truth specifications, we evaluate whether they can serve as a reliable behavioral reference for aggregate adequacy analysis. We assess relation-level coverage, gap characteristics, complementarity with execution-based coverage, and practical relevance.

\subsection{Research Questions}\label{sec:eva:rqs}
Our evaluation is organized around four research questions.

\textbf{RQ1 (MR-based Behavioral Coverage):} To what extent do existing UI component test suites exercise (\textit{Touch}) and explicitly validate (\textit{Cover}) the inferred behavioral relations?

\textbf{RQ2 (Behavioral Gaps):} Which behavioral relation types are most often left \textit{Untouched} or \textit{Weak}, and what do these gaps suggest about existing UI component tests?

\textbf{RQ3 (Complementary Adequacy Signals):} Does MR coverage provide adequacy information beyond statement and branch coverage?

\textbf{RQ4 (Practical Relevance):} Are reported issue descriptions represented in the inferred MR space, and do MR-based labels provide supporting evidence for practical behavioral validation?

\subsection{Dataset}\label{sec:dataset}
We evaluate the framework on four widely used open-source UI component libraries: \textit{ant-design}, \textit{element-plus}, \textit{material-ui}, and \textit{base-ui}. These libraries represent mature, production-grade UI ecosystems spanning both React and Vue, with diverse implementation strategies, documentation styles, and testing practices. Together, they provide a diverse benchmark for evaluating behavioral validation across modern component libraries. 
For each component, we collect its implementation, public documentation/API specification, and test suite. These artifacts provide the inputs required to construct component profiles, infer behavioral MRs, and align inferred relations with existing tests.

Table~\ref{tab:dataset} summarizes the dataset. In total, it contains 214 UI components distributed across five common categories: Inputs, Data Display, Navigation, Feedback, and Layout. The category labels follow the official classifications used by the libraries studied and are used only for descriptive analyses. All MR inference, alignment, and MR coverage measurements are performed at the individual component level.

\begin{table}[h]
\centering
\caption{Studied libraries, components, and their classification. \#Comp. denotes the number of components. $LOC$ denotes average lines of code per component.}
\label{tab:dataset}
\small
\setlength{\tabcolsep}{2pt}
\begin{tabular}{lrrrrrr}
\toprule
&  & & \textbf{ } & & \multicolumn{2}{c}{\textbf{Comp.} ($LOC$)} \\
\textbf{Library} & \textbf{Stars} & \textbf{Commits} &  \textbf{Issues$^\dagger$} & \textbf{\#Comp.} &  \textbf{Src} & \textbf{Test}  \\
\midrule
ant-design\footnotemark[1]   & 98.2k & 32{,}243 & 1{,}224 & 63 & 653   & 393  \\
element-plus\footnotemark[2] & 27.5k &  7{,}636 &   891   & 71 & 736   & 523  \\
material-ui\footnotemark[3] & 98.4k & 28{,}304 & 1{,}399 & 47 & 488   & 452  \\
base-ui\footnotemark[4]  &  9.8k &  4{,}362 &   316   & 33 & 1{,}213 & 1{,}643 \\
\midrule
\textbf{Total} & & & & \textbf{214}  & & \\
\bottomrule
\multicolumn{7}{l}{\footnotesize $^\dagger$Repository and open issues were collected as of 2026-05-10.}
\end{tabular}
\vspace{6pt}

\begin{tabular}{@{}lrrrrr@{}}
\toprule
\textbf{Category} & \textbf{Inputs} & \textbf{Data display} & \textbf{Navigation} & \textbf{Feedback} & \textbf{Layout} \\
\midrule
\#Comp.   & 69 & 56 & 36 & 34 & 19 \\
\bottomrule
\end{tabular}
\end{table}
\footnotetext[1]{\url{https://github.com/ant-design/ant-design}}
\footnotetext[2]{\url{https://github.com/element-plus/element-plus}}
\footnotetext[3]{\url{https://github.com/mui/material-ui}}
\footnotetext[4]{\url{https://github.com/mui/base-ui}}


\subsection{Measurement Validation}\label{sec:validation}
All subsequent analyses rely on the inferred MR space and the resulting MR coverage measurements. We therefore validate both MR inference and test--MR alignment, since errors in either step can affect \textit{Touch} and \textit{Cover}. All LLM-assisted inference and semantic alignment use fixed prompts, identical procedures, and a temperature of 0.2 across components.

\paragraph{MR validity and taxonomy consistency.}
We randomly sample 420 inferred MRs (approximately 10\% of the total) using stratified sampling across taxonomy categories, libraries, and component types. Two authors independently assess whether each MR is supported by the component specification, observable through component behavior, testable in principle, non-duplicative, and correctly classified within the proposed taxonomy. Disagreements are resolved through discussion and subsequently reviewed by an experienced front-end developer. 
Manual validation showed that among the sampled MRs, 88.6\% are judged usable after manual inspection, while 6.4\% are considered invalid or unsupported, and 5.0\% are identified as duplicates. The taxonomy achieves agreement of 94.8\% at the top level and 89.5\% at the fine-grained relation level. These results provide supporting evidence that the inferred MR space and taxonomy are sufficiently reliable for the subsequent MR coverage analyses.


% \begin{table}[t]
% \centering
% \caption{Manual validation of MR validity and taxonomy consistency (420 sampled MRs).}
% \label{tab:mr_quality}
% \small
% \setlength{\tabcolsep}{4pt}
% % \renewcommand{\arraystretch}{1.08}
% \begin{tabular}{lrr}
% \toprule
% \textbf{Item} & \textbf{Count} & \textbf{\%} \\
% \midrule
% \textbf{MR validity} & & \\
% Usable & \textbf{372} & \textbf{88.6} \\
% Invalid/unsupported & 27 & 6.4 \\
% Duplicate/merged & 21 & 5.0 \\
% \midrule

% \textbf{Taxonomy consistency} & & \\
% Top-level agreement & \textbf{398} & \textbf{94.8} \\
% Fine-grained agreement & \textbf{376} & \textbf{89.5} \\
% Unclassifiable & 5 & 1.2 \\
% \bottomrule
% \end{tabular}
% \end{table}


\paragraph{Alignment Validation.}
Since MR coverage depends on test--MR alignment, we evaluate the reliability of the alignment procedure and the contribution of its two evidence channels. For each LLM configuration, we sampled 100 MR--test instances from each of four groups: \textit{Touch}-positive, \textit{Touch}-negative, \textit{Cover}-positive, and \textit{Cover}-negative. After removing duplicates, the validation sets contained 386, 385, and 387 instances for DeepSeek, Gemini, and GPT, respectively. Each sampled instance was reviewed by one author. Cases with unclear evidence were independently reviewed by a second author until consensus was reached. We did not compute an inter-rater agreement statistic because the second review was performed as an adjudication step rather than an independent annotation process. We compare deterministic-only, semantic-only, and hybrid alignment, and report precision (P), recall (R), and F1 averaged over the three LLM configurations.

Table~\ref{tab:alignment-ablation} summarizes the results. Deterministic-only alignment performs well for \textit{Touch}, but has substantially lower \textit{Cover} recall because explicit behavioral validation is harder to identify than local test actions. Semantic-only alignment achieves the highest precision but lower recall, indicating that LLM-only judgments are more conservative on this sample. The hybrid variant achieves the highest F1 for both \textit{Touch} and \textit{Cover}, with no observed false negatives in the sampled validation set. Because the validation set is stratified over predicted labels, these results should be interpreted as sample-level reliability evidence rather than population-level estimates. We therefore use hybrid alignment in the main evaluation, while treating the resulting labels as aggregate analysis signals rather than ground truth.


\begin{table}[h]
\centering
\caption{Ablation of alignment evidence across LLM configurations. Results report average precision, recall, and F1 over three LLMs on manually validated MR--test instances.}
\label{tab:alignment-ablation}
\small
\setlength{\tabcolsep}{2pt}
\begin{tabular}{lcccccc}
\toprule
& \multicolumn{3}{c}{\textbf{Touch} (\%)}& \multicolumn{3}{c}{\textbf{Cover} (\%)} \\
 \cmidrule(lr){2-4} \cmidrule(lr){5-7}
\textbf{Variant} & Precision & Recall & F1 & Precision & Recall & F1 \\
\midrule
Deterministic-only
& 92.9 & 90.4 & 91.5
& 72.6 & 50.7 & 58.9 \\
Semantic-only
& 100.0 & 73.1 & 84.2
& 94.1 & 71.8 & 81.1 \\
Hybrid
& 93.4 & 100.0 & 96.6
& 81.6 & 100.0 & 89.7 \\
\bottomrule
\end{tabular}
\end{table}



% \begin{table}[h]
% \centering
% \caption{Source attribution of system-positive Touch and Cover decisions.}
% \label{tab:alignment_source_validity}
% \small
% \setlength{\tabcolsep}{2pt}
% \renewcommand{\arraystretch}{1.1}
% \scalebox{0.94}{
% \begin{tabular}{llcc}
% \toprule
% \textbf{Decision} & \textbf{Model} & \textbf{Deterministic-supported}(\%) & \textbf{Semantic-only}(\%) \\
% \midrule
% \multirow{3}{*}{Touch}
% & DeepSeek & 96.1 & 3.9 \\
% & Gemini   & 91.1 & 8.9 \\
% & GPT & 90.5 & 9.5 \\
% \midrule
% \multirow{3}{*}{Cover}
% & DeepSeek & 51.5 & 48.5 \\
% & Gemini   & 54.1 & 45.9 \\
% & GPT & 65.8 & 34.2 \\
% \midrule
% \end{tabular}
% }
% \end{table}



% \begin{table}[h]
% \centering
% \caption{Manual validation of Touch and Cover alignment labels against human review.}
% \label{tab:alignment-precision}
% \setlength{\tabcolsep}{5pt}
% \renewcommand{\arraystretch}{1.1}
% \begin{tabular}{@{}lcc@{}}
% \toprule
% \textbf{Label} & \textbf{Positive precision} & \textbf{Negative accuracy} \\
% \midrule
% Touch & 86.0\% {\footnotesize (86/100)} & 100.0\% {\footnotesize (100/100)} \\
% Cover & 95.0\% {\footnotesize (95/100)} & 100.0\% {\footnotesize (100/100)} \\
% \bottomrule
% \end{tabular}
% \end{table}




\subsection{Execution-based Coverage Baselines}\label{sec:baselines}
We compare MR coverage with statement and branch coverage as conventional execution-based baselines. For each library, we run the original test suite using its native infrastructure and collect coverage over the same component scope used for MR analysis. At the component level, \textit{Stmt} and \textit{Branch} measure code execution, while \textit{Touch} and \textit{Cover} measure behavioral reach and behavioral validation.

\subsection{Practical Relevance Assessment}\label{sec:external-validation}
We assess external relevance from three complementary perspectives: issue-related alignment, oracle augmentation for weak relations, and detection of MR-relevant injected faults.

\paragraph{Issue-based validation.} To assess the external relevance of the inferred MR space, we examine whether behaviors described in real-world issue reports can be represented by inferred MR relation types. We collected component-related GitHub issues and pull requests from the four studied libraries. Reports unrelated to component behavior (e.g., build or infrastructure issues) were excluded through manual inspection before annotation. We then report the mapping rate and analyze how mapped issue descriptions are distributed across relation types with different average MR \textit{Cover} levels. The average MR \textit{Cover} for each relation type is computed from the DeepSeek-derived MR space and corresponding hybrid-alignment results. This analysis is intended as an external relevance assessment rather than causal evidence: issue reports are influenced by project popularity, reporting practices, and maintenance processes. We therefore do not interpret issue frequency as defect proneness, but use it to assess whether reported behaviors are represented in the inferred MR space and whether MR \textit{Cover} provides additional context for understanding their behavioral validation.

\paragraph{Weak-oracle assertion study.} To assess the practical relevance of identified weak-oracle gaps, we conduct a targeted oracle-strengthening study. We sample 200 weak-oracle relations across MR types and two libraries (\textit{ant-design} and \textit{element-plus}) with reproducible test environments. For each relation, we augment the corresponding test block with a single relation-specific assertion derived from the inferred MR. Assertion templates are instantiated using behavioral cues from the MR (e.g., placement or callback identifiers), or a conservative default assertion when no concrete cue is available. Two authors inspect failing executions and separate behavioral inconsistencies from snapshot, infrastructure, or environmental artifacts.

\paragraph{MR-relevant injected fault study.}
To assess whether MR coverage labels are associated with detecting MR-relevant behavioral faults, we conduct an MR-relevant injected fault study on 200 sampled MRs: 80 \textit{Covered}, 80 \textit{Weak}, and 40 \textit{Untouched}, stratified by MR type within each status.
For each MR, we inject one relation-specific behavioral fault into the component implementation, targeting behaviors such as disabled-state handling, ARIA mappings, state synchronization, placement behavior, or event-order semantics. We then execute the original component test suite and compare detection rates across \textit{Covered}, \textit{Weak}, and \textit{Untouched} relations.